\subsection{Japanese rings: original fields, integral closures and finiteness maps}\label{algebra-proof-066-japanese-rings-original-fields-integral-closures-and-finiteness-maps}

Source: Stacks Project authors, algebra.\AlgebraIdentifierBreak{}tex:\AlgebraIdentifierBreak{}45128--\AlgebraIdentifierBreak{}45614, original authority a044\allowbreak{}46e5\allowbreak{}7ec1\allowbreak{}fbc2\allowbreak{}52a8\allowbreak{}71af\allowbreak{}cec7\allowbreak{}752f\allowbreak{}b280\allowbreak{}7b14, file SHA256 FA8B\allowbreak{}B92E\allowbreak{}58A4\allowbreak{}F78A\allowbreak{}2BD0\allowbreak{}1B3B\allowbreak{}6A4A\allowbreak{}87DE\allowbreak{}0A0D\allowbreak{}279F\allowbreak{}5DD9\allowbreak{}0641\allowbreak{}B574\allowbreak{}DD5F\allowbreak{}BFFF\allowbreak{}A4F3. The thirteen received reports retain MC-\AlgebraIdentifierBreak{}STK-\AlgebraIdentifierBreak{}ERR-\AlgebraIdentifierBreak{}0706 through MC-\AlgebraIdentifierBreak{}STK-\AlgebraIdentifierBreak{}ERR-\AlgebraIdentifierBreak{}0714. A separate source-bound correction qualifies ``any localization'' by ``nonzero''. Complete receiving arguments below are editorial material; they do not replace the translated source. Additional paper mining and optional extensions remain deferred.

\subsubsection{Integral closures, localization and the zero-ring boundary}\label{algebra-proof-066-integral-closures-localization-and-the-zero-ring-boundary}

The definitions at 45139--45145 apply to a domain \(R\) with fraction field \(K\). N-1 requires the integral closure of \(R\) in \(K\) to be a finite \(R\)-module. N-2 requires this for its integral closure in every finite field extension \(L/K\), including \(L=K\).

The source allows zero in a multiplicative subset: 1041--1047 requires only the identity and closure under multiplication. Its fraction equivalence at 1049--1065 makes \(W^{-1}R=0\) whenever \(0\in W\), since the witness \(u=0\) equates any two fractions. Zero is not a domain: the source convention at 36--37 says it has no prime ideal, and 154 identifies domains with quotients by prime ideals. In particular, its zero ideal is not prime. Taking a field \(R=k\) and \(W=\{0,1\}\) therefore shows that ``any localization'' in 45175 includes an object outside the stated definitions of N-1 and N-2. The minimal correction is ``any nonzero localization''. The familiar intended reading quantifies over localizations which remain domains; the qualification makes that reading explicit without changing either definition or declaring the usual nonzero-localization theorem false.

For \(0\notin W\), the exact map \(W^{-1}R\to K\), \(r/w\mapsto r/w\), is injective, its image is a domain and its fraction field is the original \(K\). For \(C\), the integral closure of \(R\) in the original finite extension \(L/K\), the map \[
 W^{-1}C\longrightarrow L,\qquad c/w\longmapsto c/w
\] identifies its image with the integral closure of \(W^{-1}R\) in \(L\). Here is the coefficient calculation from 7596--7627 with all denominator factors. If \[
 c^d+\sum_{i=1}^{d}a_i c^{d-i}=0,
\] then \(c/f\) satisfies \[
 (c/f)^d+\sum_{i=1}^{d}(a_i/f^i)(c/f)^{d-i}=0.
\] Conversely, if \(z/f\) satisfies a monic equation with coefficients \(a_i/f_i\), put \(F=\prod_{j=1}^{d}f_j\). After clearing the equality in a localization there is \(f'\in W\) such that \[
 (f'Fz)^d+
 \sum_{i=1}^{d}
 f^i(f')^i\left(f_i^{i-1}\prod_{j\ne i}f_j^i\right)a_i
 (f'Fz)^{d-i}=0.
\] Thus \(f'Fz\in C\), and the original fraction is \(z/f=(f'Fz)/(ff'F)\). In the present field inclusion the clearing witness can be \(f'=1\); the displayed general formula retains it. If \(C\) is finite over \(R\), its localized generators generate \(W^{-1}C\), proving both corrected localization assertions.

For the finite-cover lemma at 45184--45199, retain the original elements \(f_1,\ldots,f_n\) generating the unit ideal and the original closure \(C\). The hypothesis ``each domain \(R_{f_i}\)'' already specifies a domain chart. Alternatively, any zero chart can be removed from this cover because \(R\) is a domain and its defining element must be zero. Choose finitely many elements of \(C\) whose localizations generate each \(C_{f_i}\), clearing the denominators of the chosen module generators. Their joint \(R\)-span \(N\) is finite and \((C/N)_{f_i}=0\). For each \(v\in C/N\), choose integers \(e_i\geq1\) with \(f_i^{e_i}v=0\). The ideal generated by these powers has radical the unit ideal, so there are \(b_i\in R\) with \(1=\sum_i b_i f_i^{e_i}\). Consequently \(v=0\), and \(C=N\). The later uses of localization in this section invert nonzero elements or complements of prime ideals; none uses the excluded zero ring.

For use throughout the section, every algebraic \(\alpha\in L\) has a nonzero multiplier \(g\in R\) with \(g\alpha\) integral. Write its original monic polynomial as \[
 P(T)=T^d+\sum_{i=0}^{d-1}a_iT^i\in K[T]
\] and choose \(g\ne0\) with every \(ga_i\in R\), for example the product of the actual coefficient denominators. Keep \(\beta=g\alpha\) alongside \(\alpha\). Its equation is \[
 g^dP(T/g)=T^d+\sum_{i=0}^{d-1}g^{d-i}a_iT^i\in R[T],
 \qquad \alpha=g^{-1}\beta.
\] This also proves that the fraction field of the integral closure in \(L\) is exactly \(L\).

\subsubsection{Finite-variable descent and finite extensions of domains}\label{algebra-proof-066-finite-variable-descent-and-finite-extensions-of-domains}

We verify the infinite polynomial example at 45153--45170 without substituting a different extension. Let \(K=k(x_1,x_2,\ldots)\), and choose an actual \(K\)-basis \(e_1=1,e_2,\ldots,e_d\) of the original finite \(L/K\). Its finitely many multiplication-table coefficients lie in \(K_0=k(x_1,\ldots,x_n)\) for some \(n\). The \(K_0\)-span \(L_0\) of these same basis elements inside \(L\) is closed under multiplication and contains 1. It is a finite-dimensional domain over \(K_0\); multiplication by any nonzero element is injective and hence surjective, so \(L_0\) is a field. The map \[
 K\otimes_{K_0}L_0\longrightarrow L,\qquad a\otimes b\longmapsto ab
\] is an isomorphism because it sends the displayed basis to the original basis. A polynomial relation in the tail variables with coefficients in \(L_0\), expanded in the \(e_i\), would give a relation in that \(K\)-basis, followed by a relation in the algebraically independent original variables over \(K_0\). All its coefficients vanish. Thus these tail variables are algebraically independent over \(L_0\) and \(L=L_0(x_{n+1},x_{n+2},\ldots)\).

Put \(R_0=k[x_1,\ldots,x_n]\) and let \(C_0\) be its integral closure in \(L_0\). The polynomial N-2 proof below, iterated from the field \(k\), proves that \(C_0\) is finite over \(R_0\). This is an editorial route within the present section; the source's forward citation to the ubiquity-of-Nagata proposition remains intact. The ring \(C_0[x_{n+1},x_{n+2},\ldots]\) is integral over \(R\) and has fraction field \(L\). It is normal: an integral equation and a fraction involve only finitely many variables and therefore lie in one of the finite polynomial rings over the normal domain \(C_0\), which is normal by the earlier polynomial-normality lemma. Hence it is the exact integral closure of \(R\) in \(L\). The same finite list of \(R_0\)-module generators for \(C_0\) generates it over \(R\).

For \(R=\mathbf Z[x_1,x_2,\ldots]\), take \(K_0=\mathbf Q(x_1,\ldots,x_n)\) and \(R_0=\mathbf Z[x_1,\ldots,x_n]\). The char-zero result below starts from the Noetherian normal domain \(\mathbf Z\); the same polynomial argument applies. Both infinite polynomial rings are non-Noetherian: the ideal generated by all \(x_i\) cannot be finitely generated. Any proposed finite generators use only a finite set of variables and have zero constant term; setting those variables to zero annihilates those generators but does not annihilate a remaining variable.

For the quasi-finite ascent at 45201--45232, \(R\subset S\) induces a finite extension \(K\subset L=\operatorname{Frac}(S)\): the generic fibre is a finite-type zero-dimensional domain over \(K\), hence a finite field extension. Let \(A\) be the integral closure of \(R\) in \(S\). It is an \(R\)-submodule of the integral closure of \(R\) in \(L\), which is finite by N-2; since \(R\) is Noetherian, \(A\) is finite. The earlier quasi-finite integral-closure theorem supplies the actual elements \(g_i\in A\) with \(A_{g_i}=S_{g_i}\) covering \(\operatorname{Spec}(S)\). Nonzero charts suffice. In the finite case, for a finite \(E/\operatorname{Frac}(S)\), also \(E/K\) is finite, and \[
 \operatorname{cl}_R(E)=\operatorname{cl}_S(E).
\] Indeed integrality over \(R\) implies integrality over \(S\), while integrality over the integral \(R\)-algebra \(S\) implies integrality over \(R\) by transitivity. The finite \(R\)-generators for this closure also generate it over \(S\). Thus finite extensions preserve N-2; applying this to \(A\), and then applying localization and the finite-cover argument, proves the quasi-finite assertion.

For Laurent descent at 45234--45254, retain \(C=\operatorname{cl}_R(K)\) and \(D=\operatorname{cl}_{R[z,z^{-1}]}(K(z))\). Since \(K[z,z^{-1}]\) is normal, \(D\subset K[z,z^{-1}]\), and \(C\subset D\). Write the given finite module generators as \(f_i=\sum_j a_{ij}z^j\). For \(c\in C\), express \(c=\sum_i h_i f_i\) with \(h_i=\sum_k r_{ik}z^k\) and all sums finite. The full coefficient at exponent zero is \[
 c=\sum_{i,j}r_{i,-j}a_{ij}.
\] Thus \(C\) lies in the finite \(R\)-module \(\sum_{i,j}Ra_{ij}\); Noetherianity makes \(C\) finite. We retain that original hypothesis and do not insert an optional generalization into the review.

The omitted finite-descent proof at 45256--45266 has the following exact field construction. For N-1, \(\operatorname{cl}_R(K)\) is an \(R\)-submodule of \(\operatorname{cl}_S(\operatorname{Frac}(S))\), finite over \(S\) and then over \(R\). For N-2, retain an arbitrary original finite extension \(M/K\). Choose a \(K\)-embedding \(\iota:\operatorname{Frac}(S)\to\overline M\) and put \(E=M\iota(\operatorname{Frac}(S))\). Both fields being finite over \(K\), \(E/\iota(\operatorname{Frac}(S))\) is finite. The actual inclusion \(M\to E\) sends \(\operatorname{cl}_R(M)\) into \(\operatorname{cl}_{\iota(S)}(E)\). N-2 for \(S\) makes the latter finite over \(\iota(S)\), hence over \(R\). Noetherianity makes the original submodule finite. No linear-disjointness assertion is required.

\subsubsection{The trace lattice and the original involution example}\label{algebra-proof-066-the-trace-lattice-and-the-original-involution-example}

For 45268--45300, let \(R\) be Noetherian and normal, \(K=\operatorname{Frac}(R)\), and \(L/K\) the original finite separable extension. Starting from a \(K\)-basis \(u_i\), the coefficient-clearing calculation above gives nonzero \(g_i\in R\) and integral elements \(x_i=g_i u_i\), retaining the inverse equations \(u_i=g_i^{-1}x_i\). These still form a \(K\)-basis.

The pairing is exactly \(\langle v,w\rangle=\operatorname{Tr}_{L/K}(vw)\). The nondegeneracy theorem is fields.\AlgebraIdentifierBreak{}tex:\AlgebraIdentifierBreak{}2527--\AlgebraIdentifierBreak{}2583. Its separable conclusion can also be read on the exact scalar-extension map \[
 \overline K\otimes_K L\longrightarrow\prod_{\sigma:L\hookrightarrow\overline K}\overline K,
 \qquad a\otimes b\longmapsto(a\sigma(b))_\sigma.
\] It is an isomorphism: adjoin finitely many separable generators successively; at each step the defining polynomial factors into distinct linear factors over \(\overline K\), and the Chinese remainder evaluation maps are inverse isomorphisms. In these coordinates multiplication by \(b\) is diagonal with entries \(\sigma(b)\); its trace is their sum. For the matrix \(B=(\sigma(x_i))_{\sigma,i}\), the basis comparison makes \(B\) invertible, and the trace matrix is \(G=B^{\mathsf t}B\), so \(\det(G)=\det(B)^2\ne0\).

If \(z\) is integral over \(R\), the coefficients of its minimal polynomial \(P_z(T)=T^d+c_{d-1}T^{d-1}+\cdots+c_0\) belong to \(R\) by the normal-domain minimal-polynomial lemma. The exact trace formula from fields.\AlgebraIdentifierBreak{}tex:\AlgebraIdentifierBreak{}2425--\AlgebraIdentifierBreak{}2441 is \[
 \operatorname{Tr}_{L/K}(z)=-[L:K(z)]c_{d-1}.
\] The sign and integer factor remain, including when the integer becomes zero in the field. Thus the trace lies in \(R\). Put \[
 v_j=\sum_l(G^{-1})_{lj}x_l,\qquad
 M=\{y\in L:\operatorname{Tr}_{L/K}(x_i y)\in R\text{ for every }i\}.
\] The exact inverse maps are \[
 M\longrightarrow R^{\oplus n},\quad y\longmapsto(\operatorname{Tr}(x_i y))_i,
 \qquad (r_j)_j\longmapsto\sum_j r_jv_j.
\] Matrix multiplication gives \(\operatorname{Tr}(x_i v_j)=\delta_{ij}\), proving both composites are identities. For \(S=\operatorname{cl}_R(L)\), each \(x_i y\) is integral, so \(S\subset M\). As \(R\) is Noetherian, \(S\) is finite.

In 45302--45314 retain \(A=\mathbf C[x_1,x_2,\ldots]\), the involution \(\sigma(x_i)=-x_i\), \(L=\operatorname{Frac}(A)\), and \(R=A^{\{1,\sigma\}}\). A polynomial is invariant precisely when every monomial with nonzero coefficient has even total degree, since the monomials are a basis and \(2\ne0\). Every even-degree monomial is a product of pairs \(x_i x_j\), proving the stated generators.

The fraction field is exactly \(K=L^{\{1,\sigma\}}\). For an invariant fraction \(a=u/v\), the original nonzero denominator product \(v\sigma(v)\) belongs to \(R\), and \[
 a=\frac{u\sigma(v)}{v\sigma(v)},\qquad
 u\sigma(v)=a\,v\sigma(v)\in A\cap L^{\{1,\sigma\}}=R.
\] This proves the required surjection from \(\operatorname{Frac}(R)\) onto the invariant field; the reverse inclusion is immediate. Since \(A\) is normal by the finite-variable argument, an element of \(K\) integral over \(R\) lies in \(A\) and remains invariant, hence lies in \(R\). Thus \(R\) is normal. Each \(x_i\) satisfies \(T^2-x_i^2\in R[T]\), so \(A\) is integral over \(R\); its normality and fraction field \(L\) identify it with \(\operatorname{cl}_R(L)\).

Moreover \(L=K(x_1)\), since \(x_i=(x_1x_i)/x_1\), and \(x_1\notin K\) because \(\sigma(x_1)=-x_1\ne x_1\). Thus \([L:K]=2\), and it is separable. Let \(R_+\) be the ideal of positive-degree invariants. Its extension \(R_+A\) consists exactly of polynomials with monomials of total degree at least two. Hence \[
 A/R_+A=\mathbf C\oplus\bigoplus_{i\geq1}\mathbf C x_i
\] as a \(\mathbf C=R/R_+\)-vector space, which is infinite-dimensional. Therefore \(A\) cannot be a finite \(R\)-module. Also \(R\) is non-Noetherian: the degree-two part of \(R_+\) has the infinite independent family \(x_i x_j\) with \(i\leq j\); a finite list of ideal generators could span only finitely many degree-two components, since every multiplier of positive degree raises degree to at least four. MC-\AlgebraIdentifierBreak{}STK-\AlgebraIdentifierBreak{}ERR-\AlgebraIdentifierBreak{}0706 repairs only the apposition; the original counterexample is valid.

\subsubsection{The derivation and every original denominator factor}\label{algebra-proof-066-the-derivation-and-every-original-denominator-factor}

For 45321--45367 let \(D:R\to R\) be the original derivation. Its unique extension to \(K\) is \[
 D(r/s)=\frac{sD(r)-rD(s)}{s^2},\qquad s\ne0.
\] Uniqueness follows by applying the product rule to \(s(r/s)=r\). If \(r/s=r'/s'\), differentiating \(rs'=r's\) and multiplying by \(ss'\) gives equality of these displayed fractions; addition and multiplication follow from the same product-rule expansion. This proves the typed extension specified by MC-\AlgebraIdentifierBreak{}STK-\AlgebraIdentifierBreak{}ERR-\AlgebraIdentifierBreak{}0707.

Retain the original \(a\in K\), \(\delta_a=D(a)\ne0\), and \(L=K[x]/(x^p-a)\). Since a derivation kills \(p\)th powers, \(a\notin K^p\). The polynomial \(T^p-a\) is irreducible: the minimal polynomial of its unique root in an algebraic closure is a purely inseparable polynomial of degree dividing \(p\); if its degree were one the root would belong to \(K\). Thus \(L\) is a field, with basis \(1,x,\ldots,x^{p-1}\).

Choose an actual nonzero \(f\in R\) with \(b=f^pa\in R\); if \(a=c/d\), take \(f=d\) and \(b=d^{p-1}c\). Keep \(\beta=fx\), \(b\), \(a\), and \(x\) simultaneously. The comparison \[
 K[T]/(T^p-b)\longrightarrow L,\quad T\longmapsto fx,
 \qquad x\longmapsto f^{-1}T
\] is an isomorphism: the two relations are related by \((fx)^p-b=f^p(x^p-a)\), and the inverse uses \(f^{-p}(T^p-b)\). The derivative is \[
 \delta_b=D(b)=p f^{p-1}D(f)a+f^pD(a)=f^p\delta_a\in R\setminus\{0\}.
\] The first summand vanishes because the field has characteristic \(p\), as displayed; its contribution has not been absorbed into the notation.

We prove by induction on \(0\leq i<p\) that if \(y=\sum_{j=0}^{i}c_j\beta^j\), \(c_j\in K\), is integral over \(R\), then \(\delta_b^i c_j\in R\) for all \(j\). For \(i=0\), normality gives \(c_0\in R\). For \(i>0\), normality first gives \[
 y^p=\sum_{j=0}^{i}c_j^p b^j\in R,
 \qquad
 D(y^p)=\delta_b\sum_{j=1}^{i}j c_j^p b^{j-1}\in R.
\] Put \(z=\delta_b\sum_{j=1}^{i}j c_j\beta^{j-1}\). Its full \(p\)th power is \[
 z^p=\delta_b^p\sum_{j=1}^{i}j^p c_j^p b^{j-1}
     =\delta_b^{p-1}D(y^p)\in R,
\] because \(j^p=j\) in the prime field. Thus \(z\) is integral, and the induction hypothesis for degree \(i-1\) gives \(\delta_b^i j c_j\in R\). The integers \(1,\ldots,i\) are nonzero prime-field elements and hence units of \(R\), so \(\delta_b^i c_j\in R\) for \(j>0\). Finally \[
 \delta_b^i c_0=\delta_b^i y-\sum_{j=1}^{i}\delta_b^i c_j\beta^j
\] is integral over \(R\) and belongs to \(K\); normality gives membership in \(R\). This is exactly the missing invocation made explicit by MC-\AlgebraIdentifierBreak{}STK-\AlgebraIdentifierBreak{}ERR-\AlgebraIdentifierBreak{}0708; integrality alone would not imply membership.

The entire integral closure consequently lies in the original finite lattice \[
 \sum_{j=0}^{p-1}R\,\delta_b^{-(p-1)}\beta^j
 =\sum_{j=0}^{p-1}R\,\frac{(fx)^j}{(f^pD(a))^{p-1}}\subset L.
\] For the original coordinates \(y=\sum_j a_jx^j\), the exact comparison is \(c_j=f^{-j}a_j\). Thus no original denominator factor has been lost by the auxiliary integral presentation. The closure is an \(R\)-submodule of this finite module, so is finite by Noetherianity.

\subsubsection{Purely inseparable reduction and polynomial N-2}\label{algebra-proof-066-purely-inseparable-reduction-and-polynomial-n-2}

For 45369--45408, N-2 implies N-1 by taking \(L=K\). In characteristic zero, if N-1 holds, \(R'=\operatorname{cl}_R(K)\) is a finite Noetherian normal \(R\)-algebra. Every finite \(L/K\) is separable; the trace proof gives \(\operatorname{cl}_{R'}(L)\) finite over \(R'\), hence over \(R\), and transitivity identifies it with \(\operatorname{cl}_R(L)\).

In positive characteristic, suppose all finite purely inseparable extensions have finite closures. Given the original finite \(L/K\), choose a finite normal extension \(M/K\) containing it. The field result fields.\AlgebraIdentifierBreak{}tex:\AlgebraIdentifierBreak{}3665--\AlgebraIdentifierBreak{}3681 supplies \(F=M^{\operatorname{Aut}(M/K)}\) with \(F/K\) purely inseparable and \(M/F\) Galois. In the finite case this follows directly from the finite-automorphism fixed-field theorem: \(M/F\) has group \(\operatorname{Aut}(M/K)\); every conjugate over \(K\) of an element of \(F\) lies in \(M\) and extends to an automorphism of the normal extension, so all its distinct conjugates coincide. Its irreducible polynomial therefore has only one distinct root, which says precisely that a \(p\)-power of the element lies in \(K\). Thus \(F/K\) is purely inseparable. The closure \(A=\operatorname{cl}_R(F)\) is finite, normal, Noetherian, and has fraction field \(F\) by the earlier multiplier calculation. The trace proof makes \(\operatorname{cl}_A(M)=\operatorname{cl}_R(M)\) finite. The original \(\operatorname{cl}_R(L)=L\cap\operatorname{cl}_R(M)\) is finite as an \(R\)-submodule. The converse is contained in the definition of N-2.

For polynomial N-1 at 45410--45440, \(R'[x]\), with \(R'=\operatorname{cl}_R(K)\), is finite over \(R[x]\), normal, and has fraction field \(K(x)\). Its integral inclusion and normality identify it exactly with the integral closure of \(R[x]\). For N-2, \(R'\) remains N-2 by finite ascent proved above; finite descent reduces the assertion to \(R'\), keeping the original comparison \(R[x]\subset R'[x]\). We describe the remaining construction for the resulting normal N-2 domain, denoting it by \(A\), with field \(F\), so the original \(R\) and \(R'\) remain identifiable. Characteristic zero follows from N-1 and the preceding equivalence.

In characteristic \(p\), retain a finite purely inseparable \(L/F(x)\). Choose its finitely many field generators \(u_j\) and a common \(q=p^e\) with \[
 u_j^q=P_j(x)/Q_j(x),\quad P_j,Q_j\in F[x],\quad Q_j\ne0.
\] In an algebraic closure containing \(L\), let \(F'/F\) adjoin the unique \(q\)th roots of every coefficient of every \(P_j,Q_j\); this is a finite purely inseparable extension. Choose \(t\) with \(t^q=x\), transcendental over \(F'\). Let \(\widetilde P_j(t)\) and \(\widetilde Q_j(t)\) retain exactly the exponents of \(P_j,Q_j\) and replace each coefficient by its chosen \(q\)th root. Then \[
 \left(\frac{\widetilde P_j(t)}{\widetilde Q_j(t)}\right)^q
 =\frac{P_j(x)}{Q_j(x)}=u_j^q.
\] The denominators remain nonzero because \(t\) is transcendental. Injectivity of the \(q\)th-power map in the ambient field gives \(u_j=\widetilde P_j(t)/\widetilde Q_j(t)\), proving the actual inclusion \(L\subset F'(t)\).

Put \(C=\operatorname{cl}_A(F')\). It is finite over \(A\), normal, and has fraction field \(F'\). The ring \(C[t]\) is normal and integral over \(A[x]\) under the original map \(x\mapsto t^q\), with fraction field \(F'(t)\). These facts identify it with the integral closure of \(A[x]\) in \(F'(t)\). If \(b_1,\ldots,b_m\) generate \(C\) over \(A\), its finite \(A[x]\)-generating list is exactly \[
 b_i t^j,\qquad 1\leq i\leq m,\quad 0\leq j<q,
\] since \(t^{aq+j}=x^a t^j\). The closure in the original field \(L\) is its intersection with this ring and therefore a finite \(A[x]\)-module, since \(A[x]\) is Noetherian. MC-\AlgebraIdentifierBreak{}STK-\AlgebraIdentifierBreak{}ERR-\AlgebraIdentifierBreak{}0709 specifies the correct named extension; MC-\AlgebraIdentifierBreak{}STK-\AlgebraIdentifierBreak{}ERR-\AlgebraIdentifierBreak{}0710 identifies the middle ring. Reuse of the source symbol \(R'\) is not an additional mathematical error.

\subsubsection{Normal loci and the exact equality of local rings}\label{algebra-proof-066-normal-loci-and-the-exact-equality-of-local-rings}

For 45442--45479 let \(R\) be a Noetherian domain and \(f\ne0\) with \(R_f\) normal. Use the already reviewed Serre criterion with its original hypotheses. A domain satisfies S1; the generic local ring is a field. Therefore a nonnormal \(R_{\mathfrak p}\) has a prime \(\mathfrak q\subset\mathfrak p\) such that either \(\operatorname{depth}R_{\mathfrak q}=1\) and \(\dim R_{\mathfrak q}\geq2\), or \(\dim R_{\mathfrak q}=1\) and \(R_{\mathfrak q}\) is not regular. MC-\AlgebraIdentifierBreak{}STK-\AlgebraIdentifierBreak{}ERR-\AlgebraIdentifierBreak{}0711 changes the connective between these alternatives to ``or'', retaining the two internal conjunctions. Since normality localizes and \(R_f\) is normal, every such \(\mathfrak q\) contains \(f\).

In the first case \(f\) is a nonzerodivisor in the maximal ideal of \(R_{\mathfrak q}\), so the depth of its quotient is zero. The localized associated-prime criterion then gives \(\mathfrak q\in\operatorname{Ass}_R(R/fR)\). Its height at least two makes it embedded: minimal primes over the nonzero principal ideal have height one by the principal ideal theorem. Conversely such an embedded associated prime has depth one in the original local domain. In the second case the quotient at the height-one prime is a nonzero zero-dimensional Noetherian local ring, hence has depth zero, again giving an associated prime. The precise finite set is \[
 E=\{\mathfrak q\in\operatorname{Ass}_R(R/fR):
 \operatorname{ht}\mathfrak q\geq2\ \text{or}\
 (\operatorname{ht}\mathfrak q=1\text{ and }R_{\mathfrak q}\text{ is not regular})\}.
\] Regular height-one associated primes are not inserted into \(E\). Serre's criterion gives one containment, and localization of normality gives the other in \[
 \operatorname{Spec}(R)\setminus U=\bigcup_{\mathfrak q\in E}V(\mathfrak q).
\] It is closed. If \(f\) is a unit, the quotient is zero, \(E\) is empty and \(U=\operatorname{Spec}(R)\), as required.

For the N-1 criterion at 45481--45532, finite closure \(C\subset K\) has actual module generators \(c_i\). Choose \(f_i\ne0\) with \(f_ic_i\in R\), and retain their full product \(f=\prod_i f_i\). Then \(R_f=C_f\), a normal ring. Local N-1 follows from nonzero localization.

Conversely assume \(R_f\) normal for \(f\ne0\) and every maximal localization N-1. Any finite intermediate \(R\subset A\subset K\) has \(A_f=R_f\), because it is integral over the normal ring \(R_f\) inside \(K\). Its nonnormal locus is closed by the preceding argument; the finite map of spectra is closed, so its image \(Z_A\subset\operatorname{Spec}(R)\) is closed. MC-\AlgebraIdentifierBreak{}STK-\AlgebraIdentifierBreak{}ERR-\AlgebraIdentifierBreak{}0712 corrects the wording assigning this image its symbol.

For a maximal \(\mathfrak m\), the finite closure \(C_{\mathfrak m}=\operatorname{cl}_{R_{\mathfrak m}}(K)\) is the localization of \(C\). Express its finitely many generators as \(a_i/s_i\) with \(a_i\in C\), \(s_i\notin\mathfrak m\). Their \(R_{\mathfrak m}\)-span is also generated by the \(a_i\). The finite ring \(A=R[a_i]\) therefore has \((R\setminus\mathfrak m)^{-1}A=C_{\mathfrak m}\), which is normal. Every prime of \(A\) above \(\mathfrak m\) consequently has a normal local ring, so \(\mathfrak m\notin Z_A\). The opens \(\operatorname{Spec}(R)\setminus Z_A\) cover all primes, since every prime lies in a maximal ideal and opens are stable under passage to smaller primes. Quasi-compactness gives finitely many \(A_i\) with \(\bigcap_i Z_{A_i}=\varnothing\). The ring \(A\subset K\) generated by all these finite integral algebras is finite over \(R\).

Here is the exact receiving step retained in MC-\AlgebraIdentifierBreak{}STK-\AlgebraIdentifierBreak{}ERR-\AlgebraIdentifierBreak{}0713. Given \(\mathfrak p'\in\operatorname{Spec}(A)\), put \(\mathfrak p=\mathfrak p'\cap R\), choose \(i\) with \(\mathfrak p\notin Z_{A_i}\), and put \(\mathfrak q_i=\mathfrak p'\cap A_i\). The ring \((A_i)_{\mathfrak q_i}\) is normal. Set \(W=A_i\setminus\mathfrak q_i\). The integral extension \(A_i\subset A\) gives \[
 (A_i)_{\mathfrak q_i}\subset W^{-1}A\subset K,
\] with the first extension integral and \(K\) the fraction field of \((A_i)_{\mathfrak q_i}\). Normality therefore gives \(W^{-1}A=(A_i)_{\mathfrak q_i}\). Under this equality \(W^{-1}\mathfrak p'\) contracts to, and thus equals, the maximal ideal of that local ring. Localizing further at \(\mathfrak p'\) leaves it unchanged and proves \[
 A_{\mathfrak p'}=(A_i)_{\mathfrak q_i}.
\] This argument first localizes both rings at the same multiplicative subset. It does not rely on the generally invalid assertion that arbitrary upper-prime localization preserves integrality over the lower local ring. All local rings of \(A\) are normal, so \(A\) is normal: an element of \(K\) integral over \(A\) belongs to each \(A_{\mathfrak p'}\), and their intersection in \(K\) is \(A\). For the last intersection assertion, a fraction outside \(A\) has a proper denominator ideal, contained in some maximal ideal, so remains outside that localization. Finally any element integral over \(R\) is integral over \(A\), so \(A=C\). This proves the original N-1 assertion.

\subsubsection{Tate finiteness with the original extension retained}\label{algebra-proof-066-tate-finiteness-with-the-original-extension-retained}

Retain the source hypotheses at 45534--45589: \(R\) is a Noetherian normal domain, \(R/xR\) is a domain and N-2, and the actual map \(R\to\varprojlim_n R/x^nR\) is an isomorphism. If \(x=0\), the N-2 hypothesis already concerns \(R\). Otherwise \(x\ne0\) and is not a unit, since the quotient is a domain. In characteristic zero the trace argument proves the conclusion. In characteristic \(p\), retain an arbitrary finite purely inseparable original \(L/K\), its closure \(S=\operatorname{cl}_R(L)\), and a power \(q=p^e\) with \(L^q\subset K\).

Choose \(y^q=x\) in an algebraic closure and retain the enlargement with distinct symbols \[
 E=L(y),\qquad T=\operatorname{cl}_R(E),\qquad S=T\cap L.
\] It is finite and purely inseparable over \(K\). The \(q\)th power of any rational expression in elements of \(L\) and \(y\) is the same rational expression with all coefficients and exponents raised to the \(q\)th power; its denominator stays nonzero, and \(L^q\subset K\), \(y^q=x\in K\) show \(E^q\subset K\). If \(T\) is finite over \(R\), its \(R\)-submodule \(S\) is finite by Noetherianity. This is the return map in MC-\AlgebraIdentifierBreak{}STK-\AlgebraIdentifierBreak{}ERR-\AlgebraIdentifierBreak{}0714; the original \(L,S\) remain present throughout the argument.

For every \(t\in T\), \(t^q\in K\) is integral over \(R\), so lies in \(R\). Set \(\mathfrak p=xR\). There is exactly one prime above it, namely \[
 Q=\{t\in T:t^q\in\mathfrak p\}.
\] The field Frobenius identities show this is an ideal, and primeness of \(\mathfrak p\) shows it is prime. Its contraction is \(\mathfrak p\), since \(r^q\in\mathfrak p\) exactly when \(r\in\mathfrak p\). Any prime over \(\mathfrak p\) contains \(t\) exactly when it contains \(t^q\), proving uniqueness. Moreover \[
 Q=yT.
\] For \(t\in Q\), write the actual equation \(t^q=xr\) with \(r\in R\). Then \((t/y)^q=r\), so \(t/y\) is integral and belongs to \(T\). Conversely \((yt)^q=xt^q\in\mathfrak p\). Division by \(y\ne0\) takes place in the unchanged enlarged field \(E\).

The nonzero principal prime \(\mathfrak p\) has height one. The normal Noetherian local ring \(R_{\mathfrak p}\) is a discrete valuation ring with maximal ideal generated by the actual \(x\). By the original Krull--Akizuki theorem, algebra.\AlgebraIdentifierBreak{}tex:\AlgebraIdentifierBreak{}29292--\AlgebraIdentifierBreak{}29309, the intermediate ring \(R_{\mathfrak p}\subset T_Q\subset E\) is Noetherian. It is normal, local, nonfield, and its maximal ideal is generated by \(y\); hence it is a discrete valuation ring. The source's finite-residue-extension theorem at 29240--29264 applies.

For completeness, its needed finiteness also has the following direct dimension proof preserving the coefficient factors. Lift any finite linearly independent residue list over \(\kappa(\mathfrak p)\) to \(b_1,\ldots,b_s\in T_Q\). Suppose \(\sum_i c_i b_i=0\) with \(c_i\in K\), not all zero. Take \(m=\min_{c_i\ne0}v_x(c_i)\in\mathbf Z\), and keep \(c_i=x^m d_i\), with \(d_i\in R_{\mathfrak p}\) and at least one unit. Then the original relation is \[
 x^m\sum_i d_i b_i=0.
\] Multiplication by the explicit inverse \(x^{-m}\in E\) gives \(\sum_i d_i b_i=0\). Reducing this equality modulo \(Q T_Q\) gives a nontrivial relation in the chosen residues, a contradiction. Hence the \(b_i\) are \(K\)-linearly independent in \(E\), so \(s\leq[E:K]\). Thus \([\kappa(Q):\kappa(\mathfrak p)]\leq[E:K]\) is finite.

The domain \(T/yT=T/Q\) embeds into \(\kappa(Q)\), is integral over \(R/xR\), and hence lies in its integral closure there. That closure is finite by the N-2 hypothesis on \(R/xR\). The quotient \(R/xR\) is Noetherian, so \(T/yT\) is finite. For each \(i\geq0\), multiplication by the actual \(y^i\) gives an isomorphism \[
 T/yT\longrightarrow y^iT/y^{i+1}T.
\] It is surjective by definition. If \(y^i t=y^{i+1}u\), cancellation of the nonzero \(y^i\) in the domain \(T\) gives \(t=yu\), proving injectivity. The \(q\) successive subquotients show that \(T/y^qT=T/xT\) is finite over \(R\).

If \(z\in\bigcap_{n\geq1}x^nT\), write \(z=x^nt_n\) for every \(n\). The full exponent identity is \[
 z^q=x^{nq}t_n^q\in x^{nq}R.
\] Separatedness of the actual completion map of \(R\) gives \(\bigcap_n x^{nq}R=0\), whence \(z^q=0\), and \(z=0\) in the field \(E\).

Apply the precise finite-module theorem at 22859--22883: completeness of the ring, separatedness of the module, and finiteness modulo the ideal suffice. No completeness assumption on \(T\) is inserted. Explicitly, let \(m_1,\ldots,m_r\in T\) lift generators modulo \(xT\). Given \(t\in T\), recursively write \[
 t-\sum_{n=0}^{N}\sum_{j=1}^{r}x^n a_{j,n}m_j\in x^{N+1}T,
 \qquad a_{j,n}\in R.
\] This follows by expressing the current remainder as \(x^N u\) and choosing coefficients for the image of \(u\) modulo \(xT\). The sums \(A_j=\sum_{n\geq0}x^n a_{j,n}\) converge in the original complete ring \(R\), with their tails in \(x^{N+1}R\). The actual finite sum \(\sum_j A_jm_j\) already belongs to \(T\); its difference from \(t\) lies in every \(x^{N+1}T\), hence is zero. Therefore \(T\) is finite over \(R\), and so is the original \(S=T\cap L\). The purely inseparable criterion proves N-2 for \(R\).

\subsubsection{Power series and receiving propagation}\label{algebra-proof-066-power-series-and-receiving-propagation}

For 45592--45608 retain the original Noetherian N-2 domain \(R\). Its normal closure \(R'\) in its fraction field is finite, normal, Noetherian, and N-2 by finite ascent. The map \(R[[x]]\subset R'[[x]]\) is finite: if \(b_1,\ldots,b_m\) generate \(R'\) over \(R\), then for an original series \(\sum_n c_nx^n\) choose \(c_n=\sum_i r_{i,n}b_i\), retaining the exact identity \[
 \sum_{n\geq0}c_nx^n=\sum_{i=1}^{m}b_i\left(\sum_{n\geq0}r_{i,n}x^n\right).
\] Noetherianity of the original power-series ring is the earlier power-series theorem.

The normality argument in algebra.\AlgebraIdentifierBreak{}tex:\AlgebraIdentifierBreak{}8149--\AlgebraIdentifierBreak{}8174 retains \(w=f/g\in\operatorname{Frac}(R'[[x]])\subset\operatorname{Frac}(R')((x))\). If \(w\) is integral, the finite module generated by its powers admits a common nonzero denominator \(h\in R'[[x]]\) with \(w^e h\in R'[[x]]\) for every \(e\geq0\). Write the original leading terms as \(w=a_nx^n+\cdots\), \(h=b_mx^m+\cdots\). The leading term of their product is \(b_m a_n^e x^{m+ne}\). Since this exponent is nonnegative for every \(e\), \(n\geq0\). Each \(b_m a_n^e\in R'\). The ring \(R'[a_n]\) is therefore an \(R'\)-submodule of \(b_m^{-1}R'\), a finite module. Noetherianity makes \(R'[a_n]\) finite, so the determinant argument makes \(a_n\) integral over \(R'\), and normality gives \(a_n\in R'\). Subtract this original term \(a_nx^n\) from \(w\); the remainder is still integral, and repeat. Every coefficient of the original Laurent expansion then lies in \(R'\), with no negative exponent, so \(w\in R'[[x]]\). The zero series is already in that ring. This proves its normality with all original leading factors retained.

The actual variable \(x\) in \(R'[[x]]\) has quotient \(R'\); the ring is complete for its powers. Tate's theorem applies to that ring and variable, yielding N-2 for \(R'[[x]]\). The finite-descent proof above then gives N-2 for the original \(R[[x]]\).

Propagation within the reviewed interval is complete. The nonzero-localization qualification is used by the finite-cover, quasi-finite-ascent and N-1-localization arguments; their specified localizations remain nonzero. The prior Serre proof feeds the exact normal-locus set and the local-equality step. The prior Krull--Akizuki and finite-over-complete-ring arguments feed Tate with the original extension and all \(q\)-power factors retained. The derivation lattice, trace lattice and finite-extension descent feed the polynomial and power-series conclusions. No new optional extension is claimed. The next source interval begins at 45615, Nagata rings; only its opening through 45644 has been read here.
